\section{Sparse regression: different thresholds for different certificates}
\label{sec:regression}

Sparse regression gives a statistical example in which root exactness,
short variable-branching trees, and large convex-piece certificates have
different thresholds. The distinction depends on the relaxation and on the
signal scaling. In particular, failure to certify the planted support does
not establish failure to certify the global optimum.

\subsection{Model and perspective relaxation}
\label{regression:model}

In this section, $n$ is the sample size, $p$ the number of features, and
$1\leq k<p$ the cardinality budget. Let $X\in\R^{n\times p}$ have independent
standard Gaussian entries, with columns $x_j$, and let
\begin{equation}\label{regression:model-eq}
 y=X\beta^*+\sigma w,\qquad w\sim N(0,I_n),\qquad w\perp X.
\end{equation}
The support $S^*=\supp\beta^*$ is deterministic and has size $k$.
The ridge parameter $\lambda>0$ is deterministic. We use the unnormalized
loss throughout:
\begin{align}
 F_\lambda(S)&=\min_{b\in\R^{|S|}}
       \{\|y-X_Sb\|^2+\lambda\|b\|^2\},
       &\OPT&=\min_{|S|\leq k}F_\lambda(S),\label{regression:integer}\\
 K&=\{z\in[0,1]^p:\mathbf1^\top z\leq k\},
       &g_\lambda(z)&=y^\top
       (I_n+\lambda^{-1}X\diag(z)X^\top)^{-1}y.
       \label{regression:perspective}
\end{align}
Multiplying both terms of the objective by $1/2$ has no effect on the
results. Dividing the entire objective by $n$ changes the ridge coefficient
to $\lambda/n$.

The Boolean relaxation of
\citet{pilanciWainwrightElGhaoui2015BooleanRelaxation} has the equivalent
perspective and dual representations below. The perspective connection
is developed by \citet{xieDeng2018BooleanRelaxation}.
\begin{align}
 g_\lambda(z)&=\min_{\beta_j=0\text{ if }z_j=0}
  \left\{\|y-X\beta\|^2+
       \lambda\sum_{z_j>0}\frac{\beta_j^2}{z_j}\right\},
       \label{regression:primal}\\
 g_\lambda(z)&=\max_{a\in\R^n}h(a,z),\qquad
 h(a,z)=2a^\top y-\|a\|^2-\lambda^{-1}
       \sum_jz_j(x_j^\top a)^2.\label{regression:dual}
\end{align}
These identities follow by completing the square in the corresponding
ridge problems. The second identity makes $g_\lambda$ convex in $z$.
At $z=\mathbf1_S$, its value is $F_\lambda(S)$. Write
\begin{equation}\label{regression:node}
 r(S_0,S_1)=\min\{g_\lambda(z):z\in K,
                         z_{S_0}=0,\ z_{S_1}=1\},\qquad R=r(\varnothing,\varnothing).
\end{equation}
The two sets of fixed coordinates are disjoint. Infeasible nodes have
value $+\infty$.

For a support $S$ of size $k$, let
\begin{equation}\label{regression:residual}
 M_S=I_n+\lambda^{-1}X_SX_S^\top,\quad
 \widehat\beta_S=(X_S^\top X_S+\lambda I)^{-1}X_S^\top y,\quad
 r_S=M_S^{-1}y=y-X_S\widehat\beta_S.
\end{equation}
The normal equations give $X_S^\top r_S=\lambda\widehat\beta_S$.
The following exact condition is the selected-support certificate in
\citet[Corollary~2]{pilanciWainwrightElGhaoui2015BooleanRelaxation}.

\begin{proposition}[Exactness at a specified support]
\label{regression:exact-condition}
For $|S|=k$, put $a_j=x_j^\top r_S$. Then
\begin{equation}\label{regression:exact-condition-eq}
 R=F_\lambda(S)\quad\Longleftrightarrow\quad
 \max_{j\notin S}|a_j|\leq\min_{i\in S}|a_i|.
\end{equation}
Strict inequality implies that $S$ is the unique optimal support and that
every single wrong fixing has bound strictly above $\OPT$.
\end{proposition}
\begin{proof}
The matrix in \eqref{regression:perspective} is positive definite, and
$\partial_jg_\lambda(\mathbf1_S)=-a_j^2/\lambda$.
The first-order condition for minimizing this convex differentiable
function over $K$ says that $\mathbf1_S$ maximizes
$\sum_j a_j^2z_j$. This is precisely \eqref{regression:exact-condition-eq}.
For strictness, evaluate \eqref{regression:dual} at $r_S$. A removal of
$i\in S$ must replace its squared correlation by a smaller one; forcing
$j\notin S$ must displace a larger selected squared correlation. Both
node bounds exceed $F_\lambda(S)$. Every other feasible support lies in
one of these nodes, which proves uniqueness.
\end{proof}

\subsection{Single-variable certificates}
\label{regression:single}

At an optimal support $S$, condition C1 means that all single wrong
fixings are certified:
\begin{equation}\label{regression:c1}
 r(\{i\},\varnothing)\geq\OPT-\varepsilon\quad(i\in S),
 \qquad r(\varnothing,\{j\})\geq\OPT-\varepsilon\quad(j\notin S).
\end{equation}
Strict C1 replaces the right sides by strict inequalities above $\OPT$.
The identity $g_\lambda(\mathbf1_T)=F_\lambda(T)$ supplies the
surviving singleton bound required by the path lemma,
\cref{lattice:path}. Thus C1 gives at most $2p+1$ nodes
for every variable-branching rule if an optimal incumbent is available
when off-path nodes are examined. Strict C1 gives the same bound under
best-bound search without an initial incumbent. Node relaxation solves,
probing work, and their running times are separate costs.
Dual sufficient tests for fixing one feature underpin the safe-screening
rules of \citet{atamturkGomez2020SafeScreening}. Condition C1 here uses
the exact single-fixing relaxation values.

The full condition is sufficient, but need not be necessary for a small
certificate. If only the removal inequalities in
\eqref{regression:c1} hold, branch on the $k$ coordinates of $S$, always
continuing through their value $1$. Each sibling has a certified removal
fixing, and the budget forces the final node to be $\mathbf1_S$. This is
a certificate with $2k+1$ nodes. Thus failure of the forced-in half of
C1 alone is no lower bound on the minimum tree size.

A useful dual certificate modifies the residual to cap its unusually
large null correlations. Fix $S$, set $a=X^\top r_S$, and choose a cap
$q>0$. Put
\begin{align}
 V&=\{j\notin S:|a_j|>q\},\qquad
 H=X_V^\top M_S^{-1}X_V,\qquad
 u=H^{-1}(a_V-q\operatorname{sign}(a_V)),\label{regression:saturation}\\
 \Delta&=M_S^{-1}X_Vu,\qquad \alpha=r_S-\Delta,\qquad
 \Gamma=u^\top Hu,\quad
 m=\min_{i\in S}|x_i^\top\alpha|,\quad
 M=\max_{j\notin S}|x_j^\top\alpha|.\label{regression:saturation-quantities}
\end{align}
The empty set $V$ is allowed and gives $\Delta=0$, $\Gamma=0$.

\begin{lemma}[Capped residual certificate]
\label{regression:capped}
If $H$ is invertible, $M\leq m$, and
$\Gamma<(m^2-M^2)/\lambda$, then $S$ is uniquely optimal and every
single wrong fixing has bound at least
\begin{equation}\label{regression:capped-bound}
 F_\lambda(S)+(m^2-M^2)/\lambda-\Gamma>F_\lambda(S).
\end{equation}
Also $0\leq F_\lambda(S)-R\leq\Gamma$.
\end{lemma}
\begin{proof}
Completing the square gives
$h(\alpha,\mathbf1_S)=F_\lambda(S)-\Gamma$ and
$\|\Delta\|^2\leq\Gamma$; the cap satisfies
$X_V^\top\alpha=q\operatorname{sign}(a_V)$.
For any node, \eqref{regression:dual} gives the lower bound
\begin{equation}\label{regression:node-dual}
 2\alpha^\top y-\|\alpha\|^2-\lambda^{-1}
 \left[\sum_{i\in S_1}c_i^2+
 \operatorname{Top}_{k-|S_1|}\{c_i^2:i\notin S_0\cup S_1\}\right],
 \qquad c=X^\top\alpha,
\end{equation}
where $\operatorname{Top}_t$ sums the largest $\min(t,N)$ entries
when $N$ entries are available.
Since all selected correlations are at least all null correlations,
a wrong fixing lowers the bracket by at least $m^2-M^2$ relative to
$\sum_{i\in S}c_i^2$. This proves
\eqref{regression:capped-bound}. The root bracket equals
$\sum_{i\in S}c_i^2$, proving the root-gap bound. Every support other
than $S$ has a wrong fixing, so its value is larger.
\end{proof}

\subsection{Two sharp planted-support thresholds}
\label{regression:easy}

For the easy regime, assume $\beta_i^*=\pm b$ on $S^*$, with arbitrary
fixed signs, and fix $b,\sigma,C_0>0$. Define
\begin{equation}\label{regression:tau}
 b_\lambda=\frac{bn}{n+\lambda},\qquad
 \omega_\lambda^2=n\sigma^2+
             \frac{kb^2\lambda^2n}{(n+\lambda)^2},\qquad
 \tau_\lambda=\frac{\lambda b_\lambda}{\omega_\lambda},\qquad
 \Lambda_\lambda=\log(p\lambda/n).
\end{equation}
The ratio $\tau_\lambda$ compares a selected correlation with the
standard deviation of a null correlation. Indeed, conditional on
$(X_{S^*},w)$, the null correlations are independent
$N(0,\|r_{S^*}\|^2)$ variables.

\begin{assumption}[Easy-regime scaling]\label{regression:easy-assumption}
As $p\to\infty$,
\begin{equation}\label{regression:easy-regime}
 \log^6p\leq n\leq p,\qquad
 1\leq k\leq C_0n/\log p,\qquad
 \sqrt n\leq\lambda\leq n/\log^2p.
\end{equation}
\end{assumption}

\begin{theorem}[Root and single-variable certification]
\label{regression:easy-thresholds}
Under \cref{regression:easy-assumption}, fix $\delta\in(0,1)$.
The following statements hold with probability tending to one.
\begin{enumerate}[label=(\alph*)]
\item If $\tau_\lambda^2\geq(2+\delta)\log p$ eventually, then
      $R=\OPT=F_\lambda(S^*)$, $S^*$ is uniquely optimal, and strict C1
      holds. If $\tau_\lambda^2\leq(2-\delta)\log p$ eventually, then
      $R<F_\lambda(S^*)$.
\item If $\tau_\lambda^2\geq(2+\delta)\Lambda_\lambda$ eventually,
      then $S^*$ is uniquely optimal and every single wrong fixing has
      bound at least $\OPT+\lambda b^2/\log p$.
\item If $\tau_\lambda^2\leq(2-\delta)\Lambda_\lambda$ eventually
      and $k\geq5000/\delta^2$, then all but at most $n/\log^2p$ null
      features $j$ satisfy
      $r(\varnothing,\{j\})<F_\lambda(S^*)-\lambda b^2/2$.
\end{enumerate}
The negative statements compare with $F_\lambda(S^*)$. They imply global
root inexactness or C1 failure at $S^*$ when $S^*$ is optimal.
\end{theorem}

The proof in \cref{app:regression:easy} controls the true coefficients
uniformly, then uses the capped residual for the positive C1 statement.
The cap changes very few null correlations above the positive threshold.
Below it, many null features jointly recover the price of consuming one
unit of the cardinality budget. A single extreme null feature need not
recover that price.

\begin{corollary}[Sample sizes and a linear-tree window]
\label{regression:window}
Within \eqref{regression:easy-regime}, the following bounds describe
certification of $S^*$, with fixed positive slack in their constants.
When the sample size is above $2k\log p$, root certification is
achievable by a deterministic ridge; below that sample size it is
impossible for each allowed deterministic ridge sequence.
Single-variable certification is achievable when the sample size is above
$2k\log(p/\sqrt n)$; below this value the comparison in part~(c) of
\cref{regression:easy-thresholds} holds for each allowed deterministic
ridge sequence, with a fixed sufficiently large $k$.

More precisely, if $k=p^{\gamma+o(1)}$, $0<\gamma<1$, and
$n=a k\log p$ with fixed $2-\gamma<a<2$, a suitable deterministic
ridge gives a unique planted optimum, an inexact root, and at most
$2p+1$ nodes for every variable-branching rule under the incumbent or
best-bound conventions above. For $a>2$, a suitable ridge certifies the
root. For $a<2-\gamma$, the single-variable planted comparison fails
for each allowed deterministic ridge sequence; this does not assert a
large tree.
\end{corollary}

The ridge choice and proof are given in
\cref{app:regression:sample-sizes}. The exact bound
$\tau_\lambda^2<n/k$ for $\sigma>0$ explains the sample-size constants.
In contrast, $\lambda=\sqrt n$ gives
$\tau_\lambda\leq b/\sigma$, whereas both logarithmic thresholds diverge.
At fixed per-entry signal-to-noise ratio, this ridge eventually fails
both planted-support tests, subject to the hypotheses for their
converses. This is an asymptotic statement. In particular,
$\log^6p\leq n\leq p$ excludes the moderate dimensions of the archived
experiments discussed in \cref{sec:experiments}; those observations do
not estimate the constants in this theorem.

\subsection{Local second-moment lifts}
\label{regression:local}

The same first-order thresholds survive specified stronger relaxations
when the design has many more columns than rows. The result concerns
unbounded, signed coefficients. Bounds on coefficients supply additional
information and are outside this comparison.
Second-moment convexification and rank-one strengthening have been
studied by \citet{dongChenLinderoth2015SparseRegression} and
\citet{atamturkGomez2019RankOneConvexification}; the comparison below
specifies the local hulls and node conventions it covers.

Let $Q=X^\top X+\lambda I$ and define
\begin{equation}\label{regression:lift-objective}
 \mathcal F(\beta,B)=\|y\|^2-2y^\top X\beta+\langle Q,B\rangle.
\end{equation}
For $T\subseteq[p]$, let
\begin{equation}\label{regression:local-hull}
 \mathcal H_T=\cl\operatorname{conv}
 \{(\zeta,b,bb^\top):\zeta\in\{0,1\}^T,
                    b_j=0\text{ whenever }\zeta_j=0\}.
\end{equation}
There is no cardinality constraint inside $\mathcal H_T$.
For an integer $s\geq1$, the projected root lift $\ell_s(z)$ minimizes
\eqref{regression:lift-objective} subject to
\begin{equation}\label{regression:local-lift}
 \begin{pmatrix}1&\beta^\top\\\beta&B\end{pmatrix}\succeq0,
 \qquad (z_T,\beta_T,B_{TT})\in\mathcal H_T\quad(|T|\leq s).
\end{equation}
The node bound $L_s(S_0,S_1)$ minimizes $\ell_s(z)$ on the face in
\eqref{regression:node}. In this definition, a zero fixing constrains
$z_j$ and $\beta_j$ but leaves second moments available:
$(0,0,t)\in\mathcal H_{\{j\}}$ for every $t\geq0$.
This convention matters for the hard lower bound below.

The singleton lift $L_1$ is the relaxation with global positive
semidefiniteness and $\beta_j^2\leq z_jB_{jj}$.
Every $L_s$ dominates the perspective relaxation and is valid for the
integer problem. Also $L_s$ dominates selected block relaxations:
if $Q=Q_0+\sum_TQ_T$, where all matrices are positive semidefinite and
each $Q_T$ is supported on at most $s$ coordinates, replacing
$\beta_T^\top Q_T\beta_T$ by its exact closed indicator epigraph hull
gives a bound no larger than $L_s$. The elementary hull proof is in
\cref{app:regression:hull}.

\begin{theorem}[Thresholds for selected local lifts]
\label{regression:lift-thresholds}
Under \cref{regression:easy-assumption}, let $s=s_p\geq1$ satisfy
$s n\log p=o(p)$. Consider any valid node relaxation
$\mathcal R(S_0,S_1)$ satisfying
\begin{equation}\label{regression:lift-sandwich}
 r(S_0,S_1)\leq\mathcal R(S_0,S_1)
 \quad\text{at every node},\qquad
 \mathcal R\leq L_s
 \quad\text{at the root and all forced-in nodes }(\varnothing,\{j\}).
\end{equation}
Then parts~(a) and (b) of \cref{regression:easy-thresholds}, including
their positive margins, hold for $\mathcal R$. Its negative C1 statement
holds with $k\geq20000/\delta^2$ and at most $2n/\log^2p$ exceptions.
Consequently the planted-support sample-size thresholds and the
linear-tree window of \cref{regression:window} also hold.
\end{theorem}

For the positive tree statement, monotonicity of $\mathcal R$ under
fixings is unnecessary: an off-path node has $\mathcal R$ at least
the monotone perspective bound of a single wrong fixing. Validity gives
path-node bounds at most $\OPT$. Column deletion at zero-fixed nodes
is therefore compatible with this theorem if
\eqref{regression:lift-sandwich} holds: its upper comparisons occur
only where no coordinate is fixed to zero.

The obstruction has a precise mechanism. A perspective point is the
mean of a random feasible support. Its coefficient covariance has trace
\begin{equation}\label{regression:fractional-excess}
 \pi(z,\beta)=\sum_{z_j>0}\beta_j^2(1/z_j-1).
\end{equation}
An exact averaged fit would pay both the ridge variance $\lambda\pi$
and the design variance. The lift can cancel the latter using columns
with $z_j=0$ while paying an additional ridge cost
$o(\lambda\pi)$. Local hulls do not detect the cancellation. The
construction and both converse proofs are complete in
\cref{app:regression:completion,app:regression:lift-easy}.

\subsection{Large certificates in pure noise and at low total signal}
\label{regression:hard}

For the hard regime, change the scaling. Let $k\to\infty$,
$k/n\to0$, $\lambda/n\to0$, and
\begin{equation}\label{regression:hard-regime}
 x_p=\frac2n\log\binom pk\longrightarrow x\in(0,x_0),
 \qquad e^{-x_0}(1+2x_0)=1,\quad x_0>0.
\end{equation}
The positive root is approximately $1.2564$.
The convex-piece certificates here branch on $z$: a node is a convex
set in $z$-space, its children cover its feasible binary points, and
its bound is the infimum of the fixed projected relaxation on that set
intersected with $K$.

\begin{theorem}[Conflict packing for perspective and pairwise hulls]
\label{regression:hard-theorem}
Under \eqref{regression:hard-regime}, consider either:
\begin{enumerate}[label=(\alph*)]
\item $y$ independent of $X$, with $y\ne0$ almost surely; or
\item the planted model \eqref{regression:model-eq}, with
      $\|\beta^*\|^2/\sigma^2\to\kappa$ and
      $\kappa<e^{-x}(1+2x)-1$.
\end{enumerate}
There exist constants $c,\eta>0$ such that, with probability tending
to one, there is a family $\mathcal C$ of $k$-subsets with
\begin{equation}\label{regression:clique-size}
 |\mathcal C|\geq\exp(c k\log(p/k))=\binom pk^{c+o(1)}.
\end{equation}
For distinct $S,T\in\mathcal C$, both
$g_\lambda((\mathbf1_S+\mathbf1_T)/2)$ and
$\ell_2((\mathbf1_S+\mathbf1_T)/2)$ are strictly below
$\OPT-\eta Y$, where $Y=\|y\|^2$ in (a) and $Y=n\sigma^2$ in (b).
Thus every convex-piece $\varepsilon$-certificate using either
projected relaxation has at least $|\mathcal C|$ leaves whenever
$0\leq\varepsilon\leq\eta Y$.

For binary variable branching with certified incumbent reductions,
the corresponding node count is at least $|\mathcal C|/(p+1)$ when
each removed binary piece is charged to its node. The $\ell_2$ assertion
uses the retained-column node convention in
\eqref{regression:local-lift}; it does not cover deletion of zero-fixed
helper columns.
\end{theorem}

Two binary points in one convex leaf bring their midpoint into the
leaf. A midpoint below the required bound therefore prevents that
leaf from containing both points. This is the conflict mechanism from
\cref{sec:lattice}; it applies to variable, split, and other convex
pieces under the stated fixed-oracle convention. The reduction count
charges at most one removed piece per binary coordinate per node.

The proof in \cref{app:regression:hard} selects highly correlated
features, packs their $k$-subsets far apart, and bounds all pairwise
union fits at once. The integer optimum receives a separate lower
bound over all supports. For pairwise hulls, a uniform design event
allows the completion of every union using its complement as helpers;
this closes the additional uniformity step required by the stronger
relaxation.

The hard theorem treats bounded \emph{total} signal-to-noise ratio.
With fixed $b/\sigma$ and $k\to\infty$, the total ratio
$kb^2/\sigma^2$ diverges, so it does not overlap the easy theorem's
fixed per-coefficient scaling. For $k=p^{\gamma+o(1)}$ and
$n=a k\log p$, \eqref{regression:hard-regime} becomes
$x=2(1-\gamma)/a$; the admissible pure-noise range is
$a>2(1-\gamma)/x_0$. The proof gives a positive exponent, not a
practical estimate of it or a general complexity lower bound for
algorithms using other relaxations.

\subsection{A richer moment matrix changes the variance price}
\label{regression:richer}

The local-lift obstruction is not universal. To illustrate its scope,
add moments $Z$ and $U$ of indicators and indicator--coefficient products.
Consider the valid relaxation with objective
\eqref{regression:lift-objective} and
\begin{equation}\label{regression:full-moment}
 Y=\begin{pmatrix}1&z^\top&\beta^\top\\
 z&Z&U\\\beta&U^\top&B\end{pmatrix}\succeq0,\qquad
 \diag Z=z,\quad\diag U=\beta,
\end{equation}
with $z\in K$, $Z\mathbf1\leq kz$, the usual binary-product inequalities
$0\leq Z_{jm}\leq\min(z_j,z_m)$,
$Z_{jm}\geq z_j+z_m-1$, and
$U_{jm}^2\leq Z_{jm}B_{mm}$ for $j\ne m$.
The last inequalities follow from Cauchy--Schwarz for any genuine
support/coefficient distribution.

\begin{proposition}[Variance price in the richer lift]
\label{regression:variance-price}
At every feasible point of \eqref{regression:full-moment}, put
$A=\supp z$ and $\nu_A=\lambda_{\min}(X_A^\top X_A)$, with
$\nu_\varnothing=0$. Then
\begin{equation}\label{regression:variance-price-eq}
 \mathcal F(\beta,B)\geq\|y-X\beta\|^2+\lambda\|\beta\|^2
                           +(\lambda+\nu_A)\pi(z,\beta).
\end{equation}
\end{proposition}
The Gram-matrix proof is given in \cref{app:regression:variance}.
It uses both the full positive semidefinite moment matrix and the
product constraints. Pairwise hulls with positive semidefiniteness
only on $(1,\beta,B)$ do not supply this argument. The proposition
does not give a sharp root threshold or a tree bound for this richer
relaxation.

\subsection{A fixed-dimensional noise-normalization check}
\label{regression:pwe}

Support recovery and relaxation exactness can differ even when the
dimension is fixed and the sample size diverges. This gives a narrow
check on the Gaussian exactness statement of
\citet[Section~3.1, Theorem~2, p.~72]{pilanciWainwrightElGhaoui2015BooleanRelaxation}.

\begin{theorem}[Fixed dimension with per-entry noise]
\label{regression:pwe-theorem}
Fix $1\leq k<d$, $S\subseteq[d]$ with $|S|=k$, and
$\beta^*\in\R^d$ with support $S$. Put
$b_{\min}=\min_{i\in S}|\beta_i^*|>0$ and fix $\sigma>0$.
For each $n$, use the Gaussian model \eqref{regression:model-eq}
with $p=d$ and $\lambda=\sqrt n$. Then $S$ is the unique optimal
support with probability tending to one, but
\begin{equation}\label{regression:pwe-limit}
 \lim_{n\to\infty}\Pbb\{R=\OPT\}
 =\left[1-2\overline\Phi(b_{\min}/\sigma)\right]^{d-k}<1,
\end{equation}
where $\overline\Phi(t)=\Pbb\{N(0,1)>t\}$.
\end{theorem}

The complete proof, including the global-optimum transfer, is in
\cref{app:regression:pwe}. In the printed model, the noise variance
per observation is $\sigma^2$ and the ridge is $\sqrt n$ for the
unnormalized loss. The displayed sample-size requirement
\begin{equation}\label{regression:pwe-condition}
 n>c_0\frac{\sigma^2+\|\beta^*\|^2}{b_{\min}^2}\log d
\end{equation}
holds eventually for every fixed $c_0>0$. Equation
\eqref{regression:pwe-limit} therefore contradicts the claimed
exactness probability $1-2e^{-c_1n}$ for every fixed $c_1>0$
under those printed hypotheses. For example, $d=2$, $k=1$,
$\beta^*=(1,0)$, and $\sigma=1$ give the limit
$1-2\overline\Phi(1)$, approximately $0.683$.

This conclusion concerns that Gaussian probability statement. The
Boolean reformulation and its deterministic exactness condition remain
the ones used above. A fixed finite coefficient-to-noise ratio, however
large, cannot make the limit in \eqref{regression:pwe-limit} equal to
one. Shrinking the per-entry noise defines a different statistical
model and requires its own theorem.
